\section{Proofs for the heterogeneous conditional-graph model}
\label{app:post-proofs}

Throughout this appendix probabilities and expectations condition on the
legitimate field $\mathcal D$ of Section~\ref{sec:post-cstat-theory}. All
quantities selected from an independent training/discovery pool are fixed in
this conditional probability space. We suppress the coordinate subscript
$\theta$ where doing so introduces no ambiguity.

\subsection{Identification and legal raw residual moments}

Finite enumeration over the supported exposure law gives
\begin{align}
 \mathbb E R_u
 &=\sum_{s:e_{us}>0}e_{us}\frac{d_{us}}{e_{us}}
       \mathbb E[Y_u-\widehat\mu_{us}\mid S_u=s,\mathcal D]
   =\sum_s d_{us}(\mu_{us}-\widehat\mu_{us}),
 \label{eq:post-proof-dr}\\
 M&=D+B^{-1}\sum_u\mathbb E R_u.
 \nonumber
\end{align}
A zero coefficient contributes zero without forming $0/0$. A nonzero
coefficient at zero exposure is unsupported. The conditional mean identity
used in~\eqref{eq:post-proof-dr} must describe the desired counterfactual
mean; a probability array by itself does not establish that identification.
The nuisance need not equal these means, and the row distributions need not
coincide.

Similarly,
\begin{equation}
 \mathbb E|R_u|^p
 =\sum_{s:e_{us}>0}|d_{us}|^p e_{us}^{1-p}
       \mathbb E[|Y_u-\widehat\mu_{us}|^p\mid S_u=s,\mathcal D]
 \leq v_u.
 \label{eq:post-proof-raw-moment}
\end{equation}
The absolute value includes the sign of a contrast. For bounded responses,
$|Y_u-\widehat\mu_{us}|\leq\max\{|L_{us}-\widehat\mu_{us}|,
|U_{us}-\widehat\mu_{us}|\}$. A midpoint fit on $[0,1]$ gives the bound
$1/2$, which every eligible method retains. For a location plus noise,
Minkowski gives
\begin{equation}
 \|Y_u-\widehat\mu_{us}\|_p
 \leq|\mu_{us}-\widehat\mu_{us}|+\|\xi_{us}\|_p.
 \label{eq:post-proof-location}
\end{equation}
A known location class and known noise norm therefore supply a legal raw
residual envelope. For a sum of primitive noises, its norm is at most the sum
of their coefficient-weighted norms. No empirical tail-scale estimate is
silently promoted to a known bound.

\subsection{Sharp clipping bias and truncated fluctuations}

\begin{lemma}[Sharp clipping-bias inequality]
\label{lem:post-bias}
For $p>1$, $\tau>0$ and every $x\in\mathbb R$,
\begin{equation}
 |x-T_\tau(x)|\leq
 \frac{(p-1)^{p-1}}{p^p}\tau^{1-p}|x|^p.
 \label{eq:post-proof-bias}
\end{equation}
The constant is sharp.
\end{lemma}
\begin{proof}
For $|x|\leq\tau$ the left side is zero. Otherwise put $z=|x|/\tau>1$.
The ratio in question is $(z-1)/z^p$, whose derivative has the sign of
$p-(p-1)z$. Its maximum is attained at $z=p/(p-1)$ and equals
$(p-1)^{p-1}/p^p$. An atom at that value achieves equality. Taking expectations
in~\eqref{eq:post-proof-bias} gives
$|\mathbb E T_\tau(R)-\mathbb E R|\leq c_pv\tau^{1-p}$.
\end{proof}

For $1<p\leq 2$, checking separately $|x|\leq\tau$ and $|x|>\tau$ gives
\begin{equation}
 T_\tau(x)^2\leq\tau^{2-p}|x|^p,
 \qquad
 \operatorname{Var}(T_\tau(R))\leq v\tau^{2-p}.
 \label{eq:post-proof-second}
\end{equation}
Moreover
$|T_\tau(R)-\mathbb E T_\tau(R)|\leq 2\tau$.
The centering is used only in the proof, not in the estimator. The second
moment inequality cannot be reused for $p>2$: $x=.1$, $\tau=1$ and $p=3$
give $.01>.001$. A centered-noise bound is also insufficient for clipping
bias: a deterministic raw residual $R=2\tau$ has zero centered moment but
clipping bias $\tau$.

\subsection{Source graphs, support restriction and the MGF}

Suppose the declared primitives are jointly independent conditional on
$\mathcal D$, and every whole row vector is measurable in its listed
primitives and $\mathcal D$. Connect every pair of rows whose primitive lists
intersect. Disjoint vertex sets with no crossing edge use disjoint primitive
collections; their generated sigma-fields are independent. This proves the
strong conditional dependency-graph property. Shared outcome noise, logging
randomness and retained response channels all enter this construction.
Pairwise covariance alone does not establish the property.

For any design-active subset, disjoint vertex sets with no induced crossing
edge also have no crossing edge in the original graph. Their independence is
therefore retained. Removing an inactive vertex never removes a primitive
shared by two retained vertices. A proper color class has no internal edges,
so its row variables are jointly independent. For more than two vertices,
this follows by separating one vertex from the remainder and applying the
graph property inductively.

For a centered random variable $X$ with $|X|\leq b$ and
$\mathbb E X^2\leq a$, expansion of the exponential gives, when
$b|\lambda|<3$,
\begin{align}
 \mathbb E e^{\lambda X}
 &\leq 1+\frac{\lambda^2a}{2}
          \sum_{j=0}^\infty\left(\frac{b|\lambda|}{3}\right)^j,
 &\log\mathbb E e^{\lambda X}
 &\leq\frac{\lambda^2a}{2(1-b|\lambda|/3)}.
 \label{eq:post-proof-bernstein-mgf}
\end{align}
Here $|\mathbb E X^k|\leq ab^{k-2}$ and
$k!\geq 2\cdot 3^{k-2}$ for $k\geq 2$ justify the series bound.

Let $X_u=T_{\tau_u}(R_u)-\mathbb E T_{\tau_u}(R_u)$,
$A=\sum_uv_u\tau_u^{2-p}$ and $\tau_{\max}=\max_u\tau_u$.
For a verified proper $\chi$-coloring $C_1,\ldots,C_\chi$, H\"older's
inequality followed by within-color independence yields
\begin{align}
 \mathbb E e^{\lambda\sum_uX_u}
 &\leq\prod_{j=1}^\chi
    \left(\mathbb E e^{\chi\lambda\sum_{u\in C_j}X_u}\right)^{1/\chi},
 \nonumber\\
 \log\mathbb E e^{\lambda\sum_uX_u}
 &\leq\frac{\chi\lambda^2A}
              {2(1-(2/3)\chi\tau_{\max}|\lambda|)}.
 \label{eq:post-proof-graph-mgf}
\end{align}
Different color sizes and different row means do not alter this argument.
There is no decomposition into independent macro-block observations.

Write $V=\chi A$ and $c=(2/3)\chi\tau_{\max}$. The sub-gamma bound
$\log\mathbb E e^{\lambda\sum X_u}\leq
V\lambda^2/(2(1-c|\lambda|))$ implies
\begin{equation}
 \Pr\left(\left|\sum_uX_u\right|>
                \sqrt{2Vt}+ct\right)\leq 2e^{-t}.
 \label{eq:post-proof-tail}
\end{equation}
For $V>0$, substitute
$\lambda=\sqrt{2t/V}/(1+c\sqrt{2t/V})$ into Chernoff's inequality, and
apply the same calculation to $-\sum X_u$. The deterministic case is direct.
Combining~\eqref{eq:post-proof-tail} with the summed absolute bias from
Lemma~\ref{lem:post-bias}, and dividing by the original normalizer $B$, proves
the coordinate claim in Theorem~\ref{thm:post-core}. A finite union bound over
$\theta$ proves its family statement. No independence between coordinates,
methods or confidence levels is used.

\subsection{Tuning, locality, contrasts and equivalence}

Conditional on $\mathcal D$, every eligible design-only threshold or engine
choice is fixed, so the proof applies to that choice directly. It does not
require a union bound over candidates in a threshold grid. This reasoning
would not justify selecting by observed errors, observed intersected widths
or evaluation-trained fits. For example, independent Rademacher $X_1,X_2$
become equal conditional on evaluation-derived $\mathcal D=X_1X_2=1$;
conditioning is not an independence repair.

A zero coefficient produces an identically zero residual; dropping it leaves
the known direct term and denominator unchanged. In contrast, replacing $B$
by an active-row count generally changes the estimand. Restricting to the
induced graph proves the locality claim. With $d$ replaced by $\eta d$,
$v$ becomes $\eta^pv$ and $\tau$ becomes $\eta\tau$. The bias, square-root
and range terms each become $\eta$ times their previous value. This proves
Corollary~\ref{cor:post-locality}. A separate-level subtraction comparison
must allocate its total confidence budget to both level families; a paired
contrast uses its signed coefficient directly.

Any generic truncated graph mean with the same deterministic direct term,
residuals, envelopes, support, coloring, thresholds and confidence allocation
has exactly the estimator and radius in~\eqref{eq:post-estimator} and
\eqref{eq:post-radius}. Thus a matched implementation is algebraically equal,
including any common design-only bounded-branch dispatch. Coarser global graph
or global maximum-envelope versions are component ablations rather than
competitors denied available local information.

\subsection{Bounded Hoeffding, Bernstein and colorwise comparisons}

For each supported exposure, the residual lies between the endpoints of
\begin{equation}
 \frac{d_{us}}{e_{us}}[L_{us}-\widehat\mu_{us},
                       U_{us}-\widehat\mu_{us}],
 \label{eq:post-proof-actual-range}
\end{equation}
with endpoint order reversed when $d_{us}<0$. Taking the extrema over all
supported exposures gives $[\ell_u,u_u]$ and width $w_u=u_u-\ell_u$.
Zero-coefficient supported exposures contribute the value zero. Observed
sample extrema are not substituted for these known ranges.

For a variable $Z$ supported on $[a,b]$, write
$f(\lambda)=\log\mathbb E\exp(\lambda(Z-\mathbb EZ))$. Under exponential
tilting, $f''(\lambda)$ is the variance of a distribution still supported on
$[a,b]$. For any such distribution, variance is at most
$\mathbb E[(Z-(a+b)/2)^2]\leq(b-a)^2/4$. Integrating twice from
$f(0)=f'(0)=0$ proves Hoeffding's MGF bound
$f(\lambda)\leq\lambda^2(b-a)^2/8$ for either sign of $\lambda$.
Applying this rowwise and using the same coloring/H\"older step
bounds the summed log MGF by $\chi\lambda^2\sum_uw_u^2/8$.
A two-tail Chernoff argument gives~\eqref{eq:post-hoeffding}.
The raw second moment is bounded by
\begin{equation}
 v_u^{(2)}=\sum_{s:e_{us}>0}\frac{d_{us}^2}{e_{us}}
            \max\{|L_{us}-\widehat\mu_{us}|,
                   |U_{us}-\widehat\mu_{us}|\}^2.
\end{equation}
Variance is also at most $w_u^2/4$, so its minimum with $v_u^{(2)}$ is valid.
The centered residual has absolute magnitude at most $w_u$. Substitution into
\eqref{eq:post-proof-bernstein-mgf}--\eqref{eq:post-proof-tail}, with scale
$c=\chi\max_uw_u/3$, proves~\eqref{eq:post-bernstein}.
Both bounded radii use the same actual nuisance-dependent support and
confidence budget as the main method.

For completeness, a colorwise independent truncation comparator is valid
with positive budgets $\delta_j$ satisfying $\sum_j\delta_j\leq\delta_\theta$.
For color $C_j$, let $A_j=\sum_{u\in C_j}v_u\tau_u^{2-p}$,
$t_j=\log(2/\delta_j)$ and $\tau_{\max,j}=\max_{u\in C_j}\tau_u$.
By an independent Bernstein argument and a union bound its total radius is
\begin{equation}
 \frac 1B\sum_j\left[
     c_p\sum_{u\in C_j}v_u\tau_u^{1-p}
     +\sqrt{2A_jt_j}+\frac 23\tau_{\max,j}t_j\right].
 \label{eq:post-proof-colorwise}
\end{equation}
Every color is retained. Equivalently, color-average estimators are recombined
with their actual color-size weights divided by $B$. Equal averaging of
unequal colors or omission of a difficult color changes the target.

In the deterministic A fixture~\eqref{eq:post-A-fixture}, every residual range
has width one and the graph is empty. Direct substitution in
\eqref{eq:post-hoeffding} gives its stated radius. On
$|\widehat M-M|\leq r_H$, the lower endpoint is at least
$M-2r_H>.2$, proving the analytic detection implication without drawing data.

\subsection{Observation thinning and a scoped lower bound}
\label{app:post-lower}

We give the explicit experiment behind Proposition~\ref{prop:post-lower}.
The subgroup has $g=N\pi$ rows arranged into $n=g/\chi$ cliques. In clique $j$,
draw independent $A_j\sim\operatorname{Bernoulli}(\epsilon)$ and
$V_j\sim\operatorname{Bernoulli}(t)$, independently between cliques. Copy them
within the $\chi$ rows. Let
\begin{equation}
 t=\min\{1,(2n\epsilon)^{-1}\},\qquad
 a=\min\{H/t,\nu t^{-1/p}\},\qquad
 \eta=at=\min\{H,\nu t^{(p-1)/p}\}.
 \label{eq:post-proof-thinning}
\end{equation}
Under sign $\sigma\in\{-1,+1\}$ define
$Y_j(0)=b$ and $Y_j(1)=b+\sigma aV_j$. Logging is independent of potential
responses. Observe the exposure and the selected response. The full logging
law consists of independent Bernoulli clique assignments, and the target joint
policy assigns action one throughout the subgroup. Outside-subgroup data have
the same distribution under both signs and carry no sign information.

Known nuisance $b$ has
$\mathbb E|Y_j(1)-b|^p=a^pt\leq\nu^p$. The normalized subgroup target is
$\theta_\sigma=b+\sigma\eta$, and its raw IPW residual moment is at most
$\epsilon^{1-p}\nu^p$. To keep conditional means in $[0,1]$, choose $b=1/2$
and $H\leq 1/2$. The response amplitudes need not be bounded: bounded conditional
means and finite moments are distinct assumptions. This construction is not a
bounded-response lower bound.

Let $P_+$ and $P_-$ denote the observable laws. They agree unless at least one
clique has $A_j=V_j=1$. On that event positive and negative spike values are
disjoint. Thus
\begin{equation}
 \operatorname{TV}(P_+,P_-)
   =1-(1-\epsilon t)^n
   \leq n\epsilon t\leq\frac 12.
 \label{eq:post-proof-tv}
\end{equation}
For any estimator $T$, integrate the triangle inequality over the common
part of the two probability measures. Since the parameter distance is
$2\eta$,
\begin{equation}
 \mathbb E_+|T-\theta_+|+\mathbb E_-|T-\theta_-|
 \geq 2\eta(1-\operatorname{TV}(P_+,P_-)).
\end{equation}
The maximum risk is at least half this sum, hence at least
$\eta(1-\operatorname{TV})\geq\eta/2$, proving
\eqref{eq:post-lower}. Multiplying targets by $\pi$ multiplies the risk bound
by $\pi$.

When $2n\epsilon\geq 1$ and the mean cap does not bind, the normalized
lower-bound scale is
\[
2^{-1-(p-1)/p}\,\nu\left(\frac{\chi}{N\pi\epsilon}\right)^{(p-1)/p}.
\]
For the raw
moment it is the same constant times
$\nu\pi^{1/p}(\chi/(N\epsilon))^{(p-1)/p}$.
Equation~\eqref{eq:post-proof-thinning} records mean saturation and small
effective-sample regimes. The integer requirements $N\pi\in\mathbb N$ and
$\chi\mid N\pi$ and the explicit observable experiment are part of the claim.
The sample-size-dependent spike distribution is a minimax construction, not a
prediction for a fixed Pareto simulation.

At exactly zero exposure support, take $A_j=0$ always and
$Y_j(1)=b\pm\min\{H,\nu\}$. The observable laws are identical while the
subgroup targets differ. The same common-measure argument gives maximum
absolute risk at least $\min\{H,\nu\}$. This is nonidentification outside
the supported core model; an IPW certificate is unavailable, and the API
refuses a nonzero target coefficient at zero exposure rather than flooring it.

There is also a simple confidence-width consequence. If a random interval $I$
has coverage at least $1-\alpha$ under each sign, change of measure gives
\begin{equation}
 P_+(\theta_+\in I,\theta_-\in I)
 \geq 1-2\alpha-\operatorname{TV}(P_+,P_-)
 \geq\frac 12-2\alpha.
\end{equation}
Consequently its diameter is at least $2\eta$ with at least the displayed
probability when positive; for a raw target replace $2\eta$ by $2\pi\eta$.
This is a non-sharp two-point consequence. It proves neither optimal
$\log(m/\delta)$ factors nor sharpness for every sparse or connected graph.

\subsection{Decision boundaries and eventwise finite-policy statements}

If a true upper-risk threshold gap is $\gamma=M-a_0>2r$, then on
$|\widehat M-M|\leq r$ the lower endpoint satisfies
$\widehat M-r\geq M-2r>a_0$. The analogous negative-gap condition gives a
safe upper endpoint. Merely assuming $\gamma>r$ is insufficient: the admissible
estimate $\widehat M=M-r$ has lower endpoint $M-2r$.
Analytic sample budgets evaluate the complete radius and this sufficient
boundary. They are not new empirical power observations or necessary
information bounds. Numerical implementations must account for any declared
additional endpoint padding in their actual decision boundary.

For Corollary~\ref{cor:post-selection}, condition first on the independent
discovery pool. Its finite selected family and all subsequent plans are fixed;
the union event of Theorem~\ref{thm:post-core} includes all coordinates used in
audit and selection. On that event every certified risk upper bound is an upper
bound on the true fixed-design risk, and every certified utility lower bound
is a lower bound on its true contrast. Exact service/capacity are deterministic.
Thus a returned policy meeting all these checks is safe and useful in their
declared sense. Any false returned-policy claim is contained in the complement
of the simultaneous event. Integrating the conditional bound proves
$\Pr(\mathrm{return}\ \mathrm{and}\ \mathrm{unsafe})\leq\delta$.
If the return probability is positive, only the general bound
$\Pr(\mathrm{unsafe}\mid\mathrm{return})\leq
\delta/\Pr(\mathrm{return})$ follows. Total abstention cannot count as useful
return; reject-all fails a positive service requirement. No uncertified
fallback is covered by the argument.

The robust-slack and oracle-feasible comparators must be distinguished.
Suppose risk intervals are centered with radii $r^R_{jk}$, and utility contrast
intervals with radii $r^U_j$. On the simultaneous event a candidate $j$ whose
true risks satisfy $R_{jk}\leq b_k-2r^R_{jk}$ has all risk upper bounds feasible.
If its true utility contrast is at least the usefulness threshold plus
$2r^U_j$ and its exact service/capacity are feasible, it enters the certified
library. Selecting a maximal utility lower endpoint then gives
\begin{equation}
 U_{\mathrm{selected}}-U_0
 \geq L^U_{\mathrm{selected}}\geq L^U_j
 \geq U_j-U_0-2r^U_j.
 \label{eq:post-proof-robust-comparator}
\end{equation}
This is the eventwise robust-margin comparator. The oracle best feasible
candidate may have zero constraint slack and fail to enter the certified
library, so~\eqref{eq:post-proof-robust-comparator} is not a universal regret
bound against that different comparator. Neither comparator is supplied to
the certifier or selector as oracle input.

Finally, under simultaneous base-policy intervals, every nonnegative weight
vector summing to one satisfies
$\sum_jw_jL_j\leq\sum_jw_jM_j\leq\sum_jw_jU_j$ on the same event. This is a
deterministic inequality for all weight vectors, including data-selected ones.
A mixture of compatible joint policies is itself a joint policy, and its fixed-
auditor affine moments equal the mixture of the base moments. For signed
affine combinations, reverse the relevant endpoints when multiplying by a
negative coefficient. No new union over mixture weights is needed.
This proves Corollary~\ref{cor:post-mixtures}; it does not identify a fresh
nonlinearly clipped estimate with the mixture of certified base estimates,
or extend to policy-dependent denominators and nonlinear fairness metrics.

\section{Eligibility of the independent strong-baseline anchors}
\label{app:post-iid-anchor}

The homogeneous independent anchor has exact logging law
$e=(3/4,1/4)$, target $q=(1/2,1/2)$, fixed nuisance $1/2$ and weights
$w_s=q_s/e_s\in\{2/3,2\}$. Its response locations are $(.35,.55)$, so its
target is $.45$, threshold $.4$ and gap $.05$. These choices precede any
anchor outcome. Each observed DR score
\[
 X=\frac 12+w_S(Y-\tfrac 12)
\]
has that target; the observed outcome $Y$ alone instead averages under the
logging law. All rows share the same distribution and use independent
sources, so the i.i.d. requirement of the Catoni anchor is substantive.

Catoni requires a centered moment envelope. In the bounded $p=2$ anchor,
$Y\in[0,1]$ gives directly
\[
 \operatorname{Var}(X)\leq
 \mathbb E\!\left[\{w_S(Y-\tfrac 12)\}^2\right]
 \leq\frac 14\sum_s e_sw_s^2.
\]
For the finite-$p$ location-plus-noise anchor, let
$\delta_s=\mu_s-1/2\in[-1/2,1/2]$ and suppose
$\|\xi_s\|_p\leq K$. Here
$K=.1(1.75/(1.75-p))^{1/p}$ and $p=1.5$, with the declared finite-grid
qualification. Since $w_s\geq 1/2$, maximization over the four location
corners yields
\[
 \left|w_s\delta_s-\tfrac 12(\delta_0+\delta_1)\right|
 \leq\tfrac 12\{|w_s-\tfrac 12|+\tfrac 12\}
 =\tfrac 12w_s.
\]
Conditional Minkowski therefore gives
\[
 \mathbb E|X-\mathbb EX|^p
 \leq(\tfrac 12+K)^p\sum_s e_sw_s^p.
\]
This equals the legal raw residual envelope in this particular balanced-target,
midpoint-fit location model. It is not a general inequality replacing a
centered moment by a raw moment. The Catoni implementation also checks its
theorem's sample-size condition and retains its deterministic root-solver
debit.

For bounded OPE, $\mathbb Ew_S=1$ and $w_S\leq 2$ are known from the design,
even if a particular log observes only the smaller weight. The author-code
one-dimensional betting instance receives this design maximum. The 2024
DR/empirical-Bernstein engine receives the same logging law, target and
nuisance, with its own stated predictable-predictor truncation and mixture
parameters. These independent anchors do not authorize applying a row-order
martingale directly to an arbitrary connected source graph.

\section{Prespecified invalid-assumption controls}
\label{app:post-negative-controls}

Four supplementary controls each use 25 complete logs with $N=32768$, one
query and the same two nominal levels. Their configurations precede their
outcomes. Each compares a valid reference under the actual model, an
intentionally invalid interval labeled heuristic, and an honest refusal when
the required premise is not supplied. The invalid output is never promoted
to a certified result. Refusal has no interval and no assessable error rate;
it is not a zero-error observation.

\paragraph{Hidden shared shock.}
Every row has response $.45+H$, where $H\in\{-.25,.25\}$ is a single common
shock. The invalid description supplies an empty graph. The true sufficient
graph is complete, although it need not be stored as a quadratic adjacency
matrix. The valid reference uses a deterministic known target range. A narrow
interval obtained by pretending the rows are independent can be wrong even
though the declared array shapes and probability rows pass validation.

\paragraph{Wrong exposure law.}
The true action-one exposure is $1/5$, generated by an exact five-way lottery;
the invalid input says $1/2$. The target puts probability .8 on action one,
the response locations are $(.2,.8)$ and the nuisance is .5. The true target is
.68. The valid reference receives the true exposure law. A probability array
being normalized does not show that it equals the logger's actual law.

\paragraph{Missing clipping bias.}
The response is $.3+100B$ with $B\sim\operatorname{Bernoulli}(1/200)$,
target equal to logging and nuisance .5. Its mean is .8. At the fixed threshold
$\tau=.01$, the expected clipped DR score is .4901, giving exact clipping
bias .3099. A legal $p=1.5$ raw moment is computed analytically from the two
atoms before sampling. The valid reference pays the complete bias bound even
when it becomes vacuous. The intentionally invalid interval retains the same
estimate but omits this debit. A realized log without an extreme score cannot
justify that omission.

\paragraph{Evaluation-trained nuisance.}
With true exposure $1/5$, target action-one probability .8 and response
locations $(.2,.8)$, the invalid fit sets both state predictions in each row
equal to its observed response. Its weighted residual is zero and its direct
term is the observed outcome, whose mean is .32 rather than the target .68.
The valid reference uses the independent fixed nuisance .5. Conditioning on
the leaked fit does not restore the required identification and independence
conditions. Source provenance can reveal a declared violation, but array
schemas cannot discover every hidden training-data dependency.

Coverage and false claims in these controls describe complete logs under
known intentionally violated premises. They are supplementary diagnostics,
with pointwise Monte Carlo uncertainty, and do not enter the primary validity
families or the historical seventeen-comparison study.
